\section{Background and related work}\label{sec:background}

\subsection{Aggregating noisy labels without a gold standard}

\citet{Dawid1979} modelled each rater by a class-conditional confusion matrix and
estimated true labels and rater accuracies jointly by the EM algorithm. The model remains
the foundation of crowdsourcing and unsupervised ensemble aggregation. Extensions add item
difficulty \citep{Whitehill2009}, Bayesian priors and hierarchical structure, and
model-based agreement coefficients \citep{Warrens2010,Warrens2013}; latent class models
with local dependence have been studied in this journal in the context of model
selection and classification \citep{Oberski2013,Marbac2016,Oberski2016}. In the LLM setting the
usual recommendation is to treat each model, or each model--prompt pair, as an annotator
and apply Dawid--Skene or a descendant \citep{Camuffo2026}. This mapping inherits the
assumption of conditional independence given the true class, which LLM configurations
violate in a structured way.

\subsection{Conditional dependence in latent class models}

The diagnostic-testing literature encountered the same violation: multiple tests applied
to the same subject share unmeasured subject characteristics, so the conditionally
independent latent class model fits poorly and biases accuracy estimates
\citep{Vacek1985}. \citet{Qu1996} introduced a subject-level Gaussian random
effect with a class-specific loading, and gave diagnostics for residual dependence.
Subsequent work compared log-linear and Gaussian random-effects formulations
\citep{Xu2019}, studied crossed subject$\times$rater random effects for ordinal ratings
without a gold standard together with robustness to a misspecified random-effects
distribution \citep{Xie2014}, and produced software \citep{Beath2017}. In all of this work
the rater is a single unit. In an LLM ensemble the rater has a known factorial structure,
model$\times$prompt$\times$run, which allows the dependence to be specified from the
design rather than assumed or estimated pairwise.

\subsection{LLM annotation: variability, measurement error and remedies}

Three strands are relevant. The first documents instability: prompt paraphrases that
preserve meaning change classification accuracy substantially \citep{Barrie2024,Chen2026},
decoding is not deterministic even at nominally deterministic settings \citep{Atil2024},
and LLM judges agree with each other more than with task-specific human annotation
\citep{Bavaresco2025,Norman2026}. The second traces the consequences downstream:
\citet{Baumann2025} replicated 37 annotation tasks with 18 models and showed that
implementation choices produce Type I, II, S and M errors; \citet{Egami2024} and
\citet{Camuffo2026} formalise the measurement-error argument, showing that annotation
error correlated with covariates yields inconsistent regression estimates regardless of
average accuracy, and that model choice can account for the bulk of annotation variance;
\citet{Messing2026} extend this to evaluation pipelines and \citet{Carlson2025} give
guidelines for management research. The third strand proposes statistical remedies.
\citet{Zhang2025} reframe LLM stochasticity as measurement error and propose a Bayesian
generalised-linear version of Dawid--Skene with covariates on the latent state; their
raters remain conditionally independent. \citet{Chen2026} take the opposite route:
because LLMs share corpora and alignment their errors may be arbitrarily dependent, so
point identification is abandoned in favour of bounds on prevalence.
Prediction-powered inference \citep{Angelopoulos2023} offers a design-based route that
requires a labelled subsample. Variance-aware protocols \citep{Camuffo2026} and
error-decomposition frameworks \citep{Xu2026,Nguyen2023} recommend prompt ensembles,
replication and noise-aware aggregation but do not supply a model in which the dependence
between configurations is estimated.

\subsection{Positioning}

The literature therefore offers independence-based aggregators known to be misspecified
for LLM ensembles, bounds that discard the information in the ensemble's design, and
protocols without an accompanying inferential model. CRE-LCM occupies the middle ground:
dependence is given the structure implied by how the annotations were generated. This
yields point identification under conditions we state, interpretable variance components
that translate into annotation-design decisions, and a natural home for the covariate and
multiclass extensions of the existing approaches.
